\subsection{The complete obstruction for \texttt{R7\_5\_3895}}
\label{sec:m5-3895-cyclic}

The second quadratic ring uses the same all-order phase inventory and
coupled rotations, together with an additional integrality check on the
fusion-norm slack.  Its Frobenius--Perron dimensions, in source order, are
\[
 (1,5+2\sqrt5,5+2\sqrt5,10+4\sqrt5,
   8+4\sqrt5,8+4\sqrt5,9+4\sqrt5),
 \qquad D=720+320\sqrt5.
\]
The formal-codegree polynomial is
\[
 (t-4)(t-16)^2(t-5)^2(t^2-1440t+6400).
\]

\begin{proposition}[Finite all-word integrality]
\label{prop:m5-all-word-integrality}
Fix commuting simple classes $A_1,\ldots,A_t$ in a pseudo-unitary
modular category, and a finite spectral domain of $s$ candidate points
containing their complete possible simultaneous eigenvalue tuples.
Suppose each point lies in $\mathcal O_K^t$ for a number field $K$;
write $\lambda_i=(\lambda_i(Y))_Y\in\mathcal O_K^s$ for the
evaluation vector of the $i$th generator.  Then the unit-coefficient functional
is integral on every integer polynomial in the $A_i$ if and only if it
is integral on an additive basis of
\[
 B=\mathbb Z[1,\lambda_1,\ldots,\lambda_t]
       \subseteq\mathcal O_K^s.
\]
Such a basis is obtained by exact additive-lattice closure under the
generators, with a certified stabilization test.  There is no assumed
maximum word degree.
\end{proposition}
\begin{proof}
For a pseudo-unitary modular category $\mathcal M$, the unit weights
in simultaneous modular diagonalization are
$d_Y^2/\operatorname{FPdim}(\mathcal M)$.  If $\mathcal M=Z(\mathcal C)$ and
$D=\operatorname{FPdim}(\mathcal C)$, simultaneous modular
diagonalization gives
\[
 \tau(P)=[\mathbf1:P(A_1,\ldots,A_t)]
      =\sum_Y\frac{z_Y^2}{D^2}P(\lambda_1(Y),\ldots,\lambda_t(Y)).
\]
This is an integer for every integer polynomial $P$, and a nonnegative
integer for each monomial.  An arbitrary integer combination of
monomials need not have nonnegative unit coefficient.
The additive group $\mathcal O_K^s$ is finitely generated, hence
Noetherian.  Starting with $1$ and the eigenvalue tuples, repeatedly
adjoin products of an additive lattice basis with each generator.
This increasing chain stabilizes; equality of successive exact lattice
bases certifies closure.  The stable lattice is exactly $B$.
Additivity of $\tau$ proves the stated equivalence.  Extra candidate
points are permitted by assigning them multiplicity zero, but omitting
an actual spectral point would invalidate the implication.
\end{proof}

For the present source, $K=\mathbb Q(\sqrt5)$,
$\mathcal O_K=\mathbb Z[\varphi]$ with
$\varphi=(1+\sqrt5)/2$, and
$D=80(5+8\varphi)$, where $5+8\varphi$ is a unit.  Thus
$W_Y=6400z_Y^2/D^2\in\mathcal O_K$.  Weighted Galois invariance of
the complete actual row orbits is checked explicitly.  Consequently
\[
 \tau(P)=\frac1{12800}\sum_Y
       \operatorname{Tr}_{K/\mathbb Q}(W_YP(\lambda(Y))).
\]
For arbitrary $h\in\mathcal O_K^s$, define
$\rho(h)=\sum_Y\operatorname{Tr}_{K/\mathbb Q}(W_Yh_Y)/12800$.
Then $\rho(12800h)$ is integral, so checking integrality on $B$ is
equivalent to checking $B+12800\mathcal O_K^s$.
It suffices to compute this latter lattice modulo
$12800\mathcal O_K^s$.  Closure now takes place in a finite abelian
group, and exact row Hermite forms certify stabilization.  The resulting
congruences are linear in the orbit multiplicities and use signed
integer slacks.  They do not incorrectly impose nonnegativity on an
arbitrary lattice-basis polynomial.

\begin{lemma}[Integral norm slack for the two remaining systems]
\label{lem:m5-3895-norm-slack}
In each of the two remaining character labellings for
$R7\_5\_3895$, the fixed instance of
Proposition~\ref{prop:m5-integral-fusion-norm} has a nonnegative
\emph{integer} slack in every categorification.
\end{lemma}
\begin{proof}
The multiplier is $x=(0,\tfrac12,\tfrac12,0,1,1,0)$.
All low chord weights and the constant term are integral.  The low
product coefficients and the quartic squared norm are integers because
they are fusion multiplicities and a sum of their squares.
For a high restriction row $v$, put
$n=v_1+v_2+2v_4+2v_5$.  The integer maximum defining $g^*$ gives
\[
 2g^*(x\cdot v)=2g^*(n/2)=\lfloor n^2/8\rfloor.
\]
There are half-integral orbit costs, so their integrality cannot be
asserted separately.  Instead, the mod-two sum of the eight all-word
congruences with stored zero-based row indices
\[
 346,347,359,362,373,382,385,390
\]
has, on every phase column, coefficient twice its high-orbit cost.
It is zero on auxiliary columns and has even right-hand side.  These
congruences follow from Proposition~\ref{prop:m5-all-word-integrality};
their moduli are respectively
$6400,320,80,80,8,8,8,8$.  Thus the total high-row contribution is
integral.  These coefficient identities are independently reconstructed
on all $3925$ columns from the canonical tensor; the norm row, indexed
$345$, is not used.  The entire gap in the norm inequality is therefore
a nonnegative integer, as required for its unscaled integer slack.
\end{proof}

\begin{proposition}[General exclusion of $R7\_5\_3895$]
\label{prop:m5-3895-cyclic-exclusion}
No complex fusion category has Grothendieck ring $R7\_5\_3895$.
\end{proposition}
\begin{proof}
The complete signed reduction in
Section~\ref{sec:m5-remaining-quadratic} has two oriented signed
quotients, both gauge equivalent to the unsigned ring.  The corresponding
chosen lifts form a tensor subcategory of the sphericalization equivalent
to the original category, so sphericality descends.  Every possible
spherical dimension character belongs to the Frobenius--Perron Galois
orbit: the other characters either vanish on a simple object or disagree
on duals.  After Galois conjugation, it is enough to exclude a
pseudo-unitary categorification.  All subsequent modular constraints
belong to the center of this rank-seven category.

Complete source enumeration gives $15$ low induction frames and $163766$
candidate high rows.  Character reciprocity leaves precisely frames $6$
and $8$.  The row transport uses the actual nonnegative row obtained by
correcting the sign of the linear Galois transport.  All eight
negative-second-indicator cases, four in each frame, have exact
covector contradictions.  For the four positive labellings of frame $6$,
the conductor-five budget has $20508$ raw allocations, $126$ after Gram
compatibility, and $32$ per labelling after the inactive lattice check.
Of these $128$ branches, $112$ have rational Farkas certificates and the
remaining $16$ have complete integer trees with $775$ leaves.

In frame $8$, the same raw budget gives $166$ Gram-compatible allocations
and $32$ for each positive labelling after the lattice check.  Another
$112$ exact Farkas certificates reduce these alternatives to allocation
indices $\{0,19,160,165\}$.  Equality of the character attached to every
restriction row gives two labelling classes, $\{0,4\}$ and $\{2,6\}$.
Both representatives $0$ and $2$ are retained; their opposite
conductor-five residue choices are not identified.  The saturated
quadratic-character criterion
(Proposition~\ref{prop:prime-quadratic-saturation}) applies at every
congruence level and gives the active budget
$\sum_{Y\ \mathrm{high}}(v_{Y4}-v_{Y5})^2=8$.

For each label representative, the complete $61$-record phase inventory
of Proposition~\ref{prop:m5-power-phase-inventory} and the coupled
rotations of Proposition~\ref{prop:m5-cyclic-multiplicities} produce a
system with $3925$ nonnegative integer variables and $665$ equations.
Its necessary conditions include the restriction Gram matrix, center
traces, low modular relations, simultaneous fusion-eigenvalue
integrality, conductor-five restrictions, all $56$ nonunit low Verlinde
triples, and the all-word unit-coefficient congruences.  The latter close
an integral order on the complete candidate spectral domain; no maximum
word-degree assumption is imposed.  The fixed fusion-norm inequality
has an integer slack by Lemma~\ref{lem:m5-3895-norm-slack}.
The final verifier recomputes every low-Verlinde triple on the filtered
domain of $145$ high orbits with caching disabled and compares the entire
payload before rebuilding these systems.

After fixing a conductor-five allocation, exact entrywise positivity
and positive semidefiniteness of the residual Gram matrix remove
impossible packets and bound every retained orbit multiplicity.  A
single packet already violating either condition has multiplicity zero;
all capacities are computed by rational arithmetic.  For the first
nonunit source object, the fourth-rotation multiplicities in eigenvalue
order $(1,-1,i,-i)$ have the form
\[
 (r,9-r,s,2-s),\qquad 0\le r\le9,\quad0\le s\le2.
\]
There are $30$ choices for each of the eight label/allocation pairs.
Exact Farkas covectors exclude $225$ of these $240$ choices.  The
remaining fifteen have $s=1$; their complete integer-chain leaf counts
are as follows.  A dash denotes a choice already excluded in the
preceding $225$-case cover.
\begin{center}
\begin{tabular}{@{}rrrr@{}}\toprule
Label case & Allocation & Leaves for $r=5$ & Leaves for $r=6$\\\midrule
0&0&114&179\\
0&19&111&211\\
0&160&39&232\\
0&165&39&146\\
2&0&---&145\\
2&19&108&121\\
2&160&90&236\\
2&165&147&208\\\bottomrule
\end{tabular}
\end{center}
Together these chains have $2111$ forced integer splits and fifteen
terminal contradictions, hence $2126$ exact leaves.

Each split is the exhaustive alternative $x_j\le k$ or $x_j\ge k+1$.
After retained bounds are represented by nonnegative integer slacks,
an exact covector is nonnegative on every matrix column and strictly
negative on the target.  It excludes one side; the final covector
excludes the last retained side.  The independent chain verifier
reconstructs the source Gram matrix, row identities, fixed packets,
fourth-rotation coordinates, packet removals and capacities, then checks
every leaf.  Numerical solver outcomes are not premises.
All remaining positive-dimension cases are covered, so no general
categorification exists.
\end{proof}

The strict read-only entry point is
\path{evidence/rank7_m5_20261020/models/sqrt5/verify_r3895_cyclic_chain.py}.
It binds the original tensor and every proof input, reconstructs the
complete finite domains and checks all branches.  Partial probe logs,
missing terminal covectors and unfrozen dependencies do not pass this
whole-chain gate.  The accompanying
\texttt{R3895\_CYCLIC\_CHAIN.md} records the complete proof inventory.
